\documentclass[11pt]{article}

\usepackage[a4paper,margin=1in]{geometry}
\usepackage[utf8]{inputenc}
\usepackage[T1]{fontenc}
\usepackage{lmodern}
\usepackage{microtype}

\usepackage{amsmath,amssymb,mathtools}
\usepackage{booktabs}
\usepackage{array}
\usepackage{enumitem}
\usepackage{xcolor}
\usepackage{listings}
\usepackage[hidelinks]{hyperref}
\usepackage{cleveref}

% --- code listings (match documentation.tex's style) ------------------------
\definecolor{codegray}{gray}{0.96}
\definecolor{codecmnt}{rgb}{0.25,0.50,0.25}
\definecolor{codekw}{rgb}{0.10,0.10,0.55}
\definecolor{codestr}{rgb}{0.55,0.10,0.10}

\lstset{
  basicstyle=\ttfamily\footnotesize,
  backgroundcolor=\color{codegray},
  keywordstyle=\color{codekw}\bfseries,
  commentstyle=\color{codecmnt}\itshape,
  stringstyle=\color{codestr},
  numbers=left, numberstyle=\tiny\color{gray},
  stepnumber=1, numbersep=6pt,
  frame=single, framerule=0pt,
  breaklines=true, breakatwhitespace=true,
  showstringspaces=false,
  captionpos=b,
  language=C,
  morekeywords={int64_t,size_t,bool,struct,scalar_field_fn,bc_fn_t,face_id_t,
                bc_kind_t,problem_t,bc_spec_t,bc_face_t,
                FACE_LOWER_X,FACE_UPPER_X,FACE_LOWER_Y,FACE_UPPER_Y,
                FACE_LOWER_Z,FACE_UPPER_Z,
                BC_DIRICHLET,BC_NEUMANN,
                DIRICHLET_HOMOG_FACE,DIRICHLET_INHOMOG_FACE,
                NEUMANN_HOMOG_FACE,NEUMANN_INHOMOG_FACE}
}

\title{Multigrid Solver -- User Tutorial:\\ Adding Your Own Problem}
\author{}
\date{}

\newcommand{\Linf}{L^{\infty}}
\newcommand{\bigO}{\mathcal{O}}

\begin{document}
\maketitle

% =============================================================================
\section{Scope and prerequisites}
% =============================================================================

This tutorial walks through adding a new Poisson problem to the
solver: a right-hand side $f$, per-face boundary data, optionally an
exact solution for error reporting, and the TOML invocation that
runs it.  The code excerpts shown in \Cref{sec:example1,sec:example2}
are cut-and-paste ready -- the only edits you make are to
\texttt{src/problem\_registry.c}.

The tutorial assumes you have already:
\begin{enumerate}
\item Built the test configuration once and run \texttt{ctest}
      successfully (see \texttt{src/CLAUDE.md} or
      \texttt{Documentation.md} \S 10).
\item Installed Python~3 with \texttt{numpy} and \texttt{h5py} (used
      by the verifiers and \texttt{scripts/make\_xdmf.py}).
\item Have ParaView or VisIt available for visual inspection (optional
      but recommended).
\end{enumerate}

The reference for the discretisation itself lives in
\texttt{doc/documentation.tex}.

% =============================================================================
\section{The model problem}\label{sec:model}
% =============================================================================

The solver discretises the 3D Poisson equation
\begin{equation}
  \Delta u(x,y,z) \;=\; f(x,y,z)
  \quad\text{on}\quad
  \Omega = [a_x, b_x] \times [a_y, b_y] \times [a_z, b_z],
  \label{eq:poisson}
\end{equation}
with per-face boundary data: each of the six faces of $\Omega$
carries either a Dirichlet condition $u = g_{\text{face}}$ or a
Neumann condition $\partial u / \partial n = q_{\text{face}}$, where
$n$ is the outward normal.

To register a new problem you supply
\begin{itemize}[topsep=0pt,itemsep=2pt]
\item one C function for $f(x, y, z)$,
\item one C function per inhomogeneous face for $g$ or $q$ (you may
      use a single function that dispatches on the face id),
\item per-face flags saying which kind of BC each face carries, and
\item optionally, a C function for the exact solution $u^\star$ so
      the driver can print $\|u_h - u^\star\|_{\Linf}$ at exit.
\end{itemize}
You drop these into a \texttt{problem\_t} literal at the bottom of
\texttt{src/problem\_registry.c} and rebuild.

% =============================================================================
\section{The user-facing API}\label{sec:api}
% =============================================================================

Three types are involved.  All three live in \texttt{src/problem.h}
and \texttt{src/bc.h} and you do not modify them -- you just
instantiate them.

\subsection{Callback signatures}

Two function-pointer typedefs cover every callback you will write:
\begin{lstlisting}
typedef double (*scalar_field_fn)(double x, double y, double z);
typedef double (*bc_fn_t)(double x, double y, double z, face_id_t face);
\end{lstlisting}
\texttt{scalar\_field\_fn} is used for $f$ and (optionally) $u^\star$:
both are functions of the three spatial coordinates only.
\texttt{bc\_fn\_t} carries an extra \texttt{face} argument so a single
function can serve every face that takes data from the same formula
(see \Cref{sec:example2} for an example).

\subsection{Per-face BC kinds}

A per-face entry is a \texttt{bc\_face\_t}:
\begin{lstlisting}
struct bc_face_t {
    bc_kind_t kind;        /* BC_DIRICHLET or BC_NEUMANN */
    bool      homogeneous; /* true => value/flux is identically 0 */
    bc_fn_t   value;       /* may be NULL when homogeneous == true */
};
\end{lstlisting}
Four convenience macros are defined at the top of
\texttt{problem\_registry.c} so the registry table reads naturally:
\begin{lstlisting}
#define DIRICHLET_HOMOG_FACE          { BC_DIRICHLET, true,  NULL }
#define DIRICHLET_INHOMOG_FACE(cb)    { BC_DIRICHLET, false, (cb) }
#define NEUMANN_HOMOG_FACE            { BC_NEUMANN,   true,  NULL }
#define NEUMANN_INHOMOG_FACE(cb)      { BC_NEUMANN,   false, (cb) }
\end{lstlisting}
The six faces are addressed by name:
\begin{center}
\begin{tabular}{lll}
\toprule
Face identifier & Outward normal & Where the formula evaluates \\
\midrule
\texttt{FACE\_LOWER\_X} & $-\hat{x}$ & $x = a_x$ \\
\texttt{FACE\_UPPER\_X} & $+\hat{x}$ & $x = b_x$ \\
\texttt{FACE\_LOWER\_Y} & $-\hat{y}$ & $y = a_y$ \\
\texttt{FACE\_UPPER\_Y} & $+\hat{y}$ & $y = b_y$ \\
\texttt{FACE\_LOWER\_Z} & $-\hat{z}$ & $z = a_z$ \\
\texttt{FACE\_UPPER\_Z} & $+\hat{z}$ & $z = b_z$ \\
\bottomrule
\end{tabular}
\end{center}

\subsection{The problem record}

Everything bundles into a \texttt{problem\_t}:
\begin{lstlisting}
struct problem_t {
    const char        *name;     /* selected by TOML [problem] name = "..." */
    struct bc_spec_t   bc;       /* 6 entries, indexed by face_id_t */
    scalar_field_fn    rhs;      /* f(x,y,z); must be non-NULL */
    scalar_field_fn    u_exact;  /* optional; NULL if no closed form */
    bool               singular; /* true => mean-zero projection. Set
                                  * this only for the all-Neumann case,
                                  * which we recommend you avoid. */
};
\end{lstlisting}
You append your new record to the \texttt{g\_problems[]} array near
the bottom of \texttt{problem\_registry.c}.  The array is
NULL-terminated -- the existing sentinel record at the very end stays
where it is; insert your record above it.

% =============================================================================
\section{Example 1 -- All-Dirichlet}\label{sec:example1}
% =============================================================================

\subsection{The problem}

We pick a manufactured solution
\begin{equation}
  u^\star(x,y,z) \;=\; \sin(\pi x)\,\sin(\pi y)\,\sin(\pi z)
  \quad\text{on}\quad [0,1]^3,
  \label{eq:ex1-u}
\end{equation}
so the right-hand side is
\begin{equation}
  f(x,y,z) \;=\; \Delta u^\star \;=\; -3\pi^2\,u^\star(x,y,z).
  \label{eq:ex1-f}
\end{equation}
Every face has $u^\star \equiv 0$ (one of $\sin(\pi\cdot 0)$ or
$\sin(\pi\cdot 1)$ is always zero on the unit cube), so all six faces
carry \emph{homogeneous} Dirichlet data and we do not need a boundary
callback at all.

\subsection{The code}

Paste the following block into \texttt{src/problem\_registry.c},
above the sentinel \texttt{\{ .name = NULL \}} record at the end of
\texttt{g\_problems[]}.  The two static helpers go above the array;
the literal goes inside it.

\begin{lstlisting}[caption={Example 1: code to add to {\tt src/problem\_registry.c}.}]
/* ============================================================
 * tutorial_example_1
 *
 *   u*(x,y,z) = sin(pi x) sin(pi y) sin(pi z)
 *   f (x,y,z) = -3 pi^2 u*
 *   homogeneous Dirichlet on every face.
 * ============================================================ */
static double tut1_rhs(double x, double y, double z)
{
    const double pi  = 3.14159265358979323846;
    const double pi2 = pi * pi;
    return -3.0 * pi2 * sin(pi * x) * sin(pi * y) * sin(pi * z);
}

static double tut1_exact(double x, double y, double z)
{
    const double pi = 3.14159265358979323846;
    return sin(pi * x) * sin(pi * y) * sin(pi * z);
}

/* ---- inside g_problems[], before the {.name=NULL} sentinel ---- */
{
    .name = "tutorial_example_1",
    .bc   = {
        .face = {
            [FACE_LOWER_X] = DIRICHLET_HOMOG_FACE,
            [FACE_UPPER_X] = DIRICHLET_HOMOG_FACE,
            [FACE_LOWER_Y] = DIRICHLET_HOMOG_FACE,
            [FACE_UPPER_Y] = DIRICHLET_HOMOG_FACE,
            [FACE_LOWER_Z] = DIRICHLET_HOMOG_FACE,
            [FACE_UPPER_Z] = DIRICHLET_HOMOG_FACE,
        },
    },
    .rhs      = tut1_rhs,
    .u_exact  = tut1_exact,
    .singular = false,
},
\end{lstlisting}

\subsection{The TOML parameter file}

\begin{lstlisting}[language={},caption={\tt tutorial1.toml}]
[grid]
nx_cells = 64
ny_cells = 64
nz_cells = 64

[solver]
multigrid = true
omega     = 1.5
n_smooth  = 2
n_iters   = 40
tol       = 1.0e-12
subcycles = 1
min_cells = 2

[problem]
name = "tutorial_example_1"

[output]
dir = "out_tut1"
\end{lstlisting}

\subsection{Build and run}

\begin{lstlisting}[language=bash]
cmake --build build -j
mpirun -np 1 build/src/driver_multigrid tutorial1.toml
\end{lstlisting}
The driver prints one progress line per V-cycle and, at exit, the
quantity
\[
  \|u_h - u^\star\|_{\Linf} \;\equiv\; \max_{ijk}\,
  | u_h(x_i, y_j, z_k) - u^\star(x_i, y_j, z_k) |.
\]
On a uniform $\bigO(h^2)$ scheme this should approximately quarter
each time you double \texttt{nx\_cells} (etc.).  Try $N = 32, 64,
128$ and confirm the ratio of successive errors is close to 4.

\subsection{Visualise}

\begin{lstlisting}[language=bash]
python3 scripts/make_xdmf.py --dir out_tut1
paraview out_tut1/multigrid.xmf
\end{lstlisting}
ParaView shows a smooth bump that peaks at $u \approx 1$ at the
domain centre and vanishes on every face -- exactly $u^\star$.

% =============================================================================
\section{Example 2 -- Mixed: three Dirichlet + three Neumann}\label{sec:example2}
% =============================================================================

\subsection{The problem}

This example uses a manufactured solution where every face carries
genuinely non-zero data, so you exercise both the Dirichlet and the
Neumann callback paths:
\begin{equation}
  u^\star(x,y,z) \;=\; e^{x+y+z}
  \quad\text{on}\quad [0,1]^3.
  \label{eq:ex2-u}
\end{equation}
Then $\Delta u^\star = 3\,u^\star$, so
\begin{equation}
  f(x,y,z) \;=\; 3\,e^{x+y+z}.
  \label{eq:ex2-f}
\end{equation}
We pick \emph{inhomogeneous Dirichlet} on the three lower faces and
\emph{inhomogeneous Neumann} on the three upper faces:
\[
\begin{array}{lcl}
  \text{lower-}x\;(x=0) & : & u = e^{y+z}, \\
  \text{lower-}y\;(y=0) & : & u = e^{x+z}, \\
  \text{lower-}z\;(z=0) & : & u = e^{x+y}, \\
  \text{upper-}x\;(x=1) & : & \partial u / \partial n = +\partial u / \partial x = e^{1+y+z}, \\
  \text{upper-}y\;(y=1) & : & \partial u / \partial n = +\partial u / \partial y = e^{x+1+z}, \\
  \text{upper-}z\;(z=1) & : & \partial u / \partial n = +\partial u / \partial z = e^{x+y+1}.
\end{array}
\]
With three Dirichlet faces present the system is non-singular -- no
mean-zero projection is needed and the V-cycle converges to the
discrete solution without any iterative-side workarounds.

\subsection{The outward-normal convention}

The single most common gotcha when adding a Neumann problem is the
sign of $q = \partial u / \partial n$:
\begin{center}
\begin{tabular}{lcl}
\toprule
Face & Outward normal & $q$ formula \\
\midrule
\texttt{FACE\_LOWER\_X} & $-\hat{x}$ & $-\,\partial u/\partial x$ \\
\texttt{FACE\_UPPER\_X} & $+\hat{x}$ & $+\,\partial u/\partial x$ \\
\texttt{FACE\_LOWER\_Y} & $-\hat{y}$ & $-\,\partial u/\partial y$ \\
\texttt{FACE\_UPPER\_Y} & $+\hat{y}$ & $+\,\partial u/\partial y$ \\
\texttt{FACE\_LOWER\_Z} & $-\hat{z}$ & $-\,\partial u/\partial z$ \\
\texttt{FACE\_UPPER\_Z} & $+\hat{z}$ & $+\,\partial u/\partial z$ \\
\bottomrule
\end{tabular}
\end{center}
A useful sanity check: if you accidentally flip the sign on one face,
the discrete error will not converge to zero as $h \to 0$.

\subsection{The code}

\begin{lstlisting}[caption={Example 2: code to add to {\tt src/problem\_registry.c}.}]
/* ============================================================
 * tutorial_example_2
 *
 *   u*(x,y,z) = exp(x + y + z)
 *   f (x,y,z) = 3 exp(x + y + z)
 *   Lower faces: inhomogeneous Dirichlet
 *   Upper faces: inhomogeneous Neumann
 *
 * For u* = e^(x+y+z) the outward-normal derivative on every upper
 * face equals u* evaluated at that face (since du/dx = du/dy =
 * du/dz = u*).  A single bc_fn_t therefore serves all three Neumann
 * faces; the face argument is ignored because the caller has
 * already pinned the appropriate coordinate to 0 or 1 before
 * invoking the callback.
 * ============================================================ */

static double tut2_rhs(double x, double y, double z)
{
    return 3.0 * exp(x + y + z);
}

static double tut2_exact(double x, double y, double z)
{
    return exp(x + y + z);
}

static double tut2_dirichlet(double x, double y, double z, face_id_t face)
{
    (void)face;             /* the caller pins x/y/z=0 on each lower face */
    return exp(x + y + z);
}

static double tut2_neumann(double x, double y, double z, face_id_t face)
{
    (void)face;             /* du/dx = du/dy = du/dz = u* for this u*    */
    return exp(x + y + z);
}

/* ---- inside g_problems[], before the {.name=NULL} sentinel ---- */
{
    .name = "tutorial_example_2",
    .bc   = {
        .face = {
            [FACE_LOWER_X] = DIRICHLET_INHOMOG_FACE(tut2_dirichlet),
            [FACE_LOWER_Y] = DIRICHLET_INHOMOG_FACE(tut2_dirichlet),
            [FACE_LOWER_Z] = DIRICHLET_INHOMOG_FACE(tut2_dirichlet),
            [FACE_UPPER_X] = NEUMANN_INHOMOG_FACE  (tut2_neumann),
            [FACE_UPPER_Y] = NEUMANN_INHOMOG_FACE  (tut2_neumann),
            [FACE_UPPER_Z] = NEUMANN_INHOMOG_FACE  (tut2_neumann),
        },
    },
    .rhs      = tut2_rhs,
    .u_exact  = tut2_exact,
    .singular = false,
},
\end{lstlisting}

A common alternative pattern, when the data on each face takes a
\emph{different} formula, is to dispatch on the \texttt{face} argument
inside one callback rather than registering a separate function per
face.  The skeleton looks like
\begin{lstlisting}
static double my_neumann(double x, double y, double z, face_id_t face)
{
    switch (face) {
    case FACE_LOWER_X: return /* formula for q at x = a_x */;
    case FACE_UPPER_X: return /* formula for q at x = b_x */;
    case FACE_LOWER_Y: return /* formula for q at y = a_y */;
    /* ... */
    default: return 0.0;
    }
}
\end{lstlisting}
Either style is fine; \texttt{tut2\_neumann} above happens to be the
same expression on every face so the dispatch is unnecessary.

\subsection{The TOML parameter file}

\begin{lstlisting}[language={},caption={\tt tutorial2.toml}]
[grid]
nx_cells = 64
ny_cells = 64
nz_cells = 64

[solver]
multigrid = true
omega     = 1.5
n_smooth  = 2
n_iters   = 40
tol       = 1.0e-12
subcycles = 1
min_cells = 2

[problem]
name = "tutorial_example_2"

[output]
dir = "out_tut2"
\end{lstlisting}

\subsection{Build, run, visualise}

\begin{lstlisting}[language=bash]
cmake --build build -j
mpirun -np 8 build/src/driver_multigrid tutorial2.toml
python3 scripts/make_xdmf.py --dir out_tut2
paraview out_tut2/multigrid.xmf
\end{lstlisting}
Expect the same rate-2 behaviour as \Cref{sec:example1}: the
$\|u_h - u^\star\|_{\Linf}$ printed at exit should drop by a factor of
roughly 4 each time you double \texttt{nx\_cells}.

% =============================================================================
\section{Verifying your solution}\label{sec:verify}
% =============================================================================

Three complementary checks:

\begin{enumerate}
\item \textbf{Closed-form error.}  When you supply \texttt{.u\_exact},
      the driver prints $\|u_h - u^\star\|_{\Linf}$ at exit.  Running
      the same TOML at $N = 32, 64, 128$ should give a ratio of
      successive errors near~4 (rate~2).  This is the strongest check
      available -- if it passes, the discretisation, BCs, and RHS are
      all consistent.

\item \textbf{Defect history.}  Each iteration prints
      $|d_h|_{\Linf} = |L_h u_h - f_h|_{\Linf}$.  It should drop
      monotonically and reach \texttt{tol} well before
      \texttt{n\_iters} is exhausted.  A flat plateau at $\sim 10^{-2}$
      almost always means the system is singular -- check that at
      least one face is Dirichlet (see \Cref{sec:pitfalls}).

\item \textbf{Visual.}  Open \texttt{multigrid.xmf} in ParaView and
      eyeball the field for the expected shape, smoothness, and
      boundary behaviour.  Cell-centred ghost slots are filtered out
      by \texttt{make\_xdmf.py}; Dirichlet vertex slots are rendered
      at the physical boundary.
\end{enumerate}

% =============================================================================
\section{Conventions and pitfalls}\label{sec:pitfalls}
% =============================================================================

\paragraph{Domain bounds.}
The default domain is $[0,1]^3$.  To use a different box, add
\texttt{x0}, \texttt{xN}, \texttt{y0}, \texttt{yN}, \texttt{z0},
\texttt{zN} keys to the \texttt{[grid]} section of the TOML.  Your
callbacks receive \emph{user-frame} coordinates, i.e. the same
$(x,y,z)$ values you used in your formulas above.  The cell-centred
internal coordinate shift on Neumann-bearing axes (see
\texttt{Documentation.md} \S2) is invisible to your callbacks.

\paragraph{Outward-normal sign for Neumann data.}
Repeated from \Cref{sec:example2} because it is the single most
common bug: $q$ is the \emph{outward}-normal derivative, so $q$ on
the lower-$x$ face is $-\partial u/\partial x$ at $x = a_x$, while
$q$ on the upper-$x$ face is $+\partial u/\partial x$ at $x = b_x$.
If your error fails to converge to zero, this is the first place to
look.

\paragraph{Avoid all-Neumann.}
The all-Neumann problem is singular: $u + c$ is a solution whenever
$u$ is.  The discrete operator inherits this null space, and the
V-cycle's mean-zero projection plus per-iteration re-injection of
the constant mode interact badly -- the defect plateaus at
$\sim 10^{-2}$ even though the discretisation itself is rate~2.
Pin at least one face to Dirichlet.  See \texttt{Plan.md} "Known
limitations" for the diagnosis.

\paragraph{Compatibility.}
With even a single Dirichlet face present, the system is
non-singular and no additional care is needed.  No discrete
compatibility condition between $f$ and $q$ needs to be verified.

\paragraph{Names and uniqueness.}
The \texttt{.name} field is what you cite from the TOML's
\texttt{[problem] name = "..."} key.  It must be unique across
\texttt{g\_problems[]}; \texttt{problem\_lookup} returns the first
match in array order.

% =============================================================================
\section{Where to look next}
% =============================================================================

\begin{description}[leftmargin=2em,style=nextline]
\item[\texttt{src/problem.h}, \texttt{src/bc.h}]
  Type definitions and the convenience macros for per-face BC entries.
\item[\texttt{src/problem\_registry.c}]
  Six built-in presets, ordered from simple to complex.  Pick the one
  closest to your new problem and copy the structure -- the existing
  Dirichlet-inhomog and mixed-inhomog presets are particularly
  instructive.
\item[\texttt{src/problem\_helpers.c}]
  The L${}^\infty$ error reduction and the mean-zero projection.  You
  do not need to touch this file when adding a new problem; it
  documents the helpers the driver calls into.
\item[\texttt{doc/documentation.tex}]
  The reference description of the discretisation and the V-cycle.
\item[\texttt{Plan.md}]
  Outstanding limitations and design notes, including the singular
  all-Neumann issue mentioned above.
\end{description}

\end{document}
